\chapter{Natural Deduction as Tactics}
\label{ch:natural-deduction}

This chapter reads Lean's core tactics directly as introduction and
elimination rules, rather than as arbitrary programming instructions
to be memorized separately from the logic they express. By the end
of the chapter, the tactic script and the natural-deduction proof
should be visibly the same object viewed two ways.

\section{Introduction Rules: \texttt{intro} and \texttt{constructor}}

Two tactics do almost all of the introduction work needed so far.
\texttt{intro} discharges a goal whose type is an implication or a
universal statement by naming its hypothesis or its bound variable
and leaving the rest of the goal to be proved with that name in
scope — this is $\to$-introduction and $\forall$-introduction,
respectively, and it is the tactic-mode counterpart of the
\texttt{fun h => \ldots} term-mode proofs used throughout
Chapter~\ref{ch:propositions}. \texttt{constructor} discharges a
goal by applying the head constructor of whatever inductive type the
goal's type is built from, splitting into one new goal per
constructor argument — for $p \land q$ that means two new goals, $p$
and $q$ (\texttt{And}-introduction), and for $\exists x, P\,x$ it
means supplying a witness and then a proof that the witness works
($\exists$-introduction).

\begin{experiment}{Building a conjunction with \texttt{constructor}}
\begin{lean}
theorem and_intro_tactic (p q : Prop) (hp : p) (hq : q) : p ∧ q := by
  constructor
  · exact hp
  · exact hq
\end{lean}
The natural-deduction rule this tactic script enacts step by step is
\[
\dfrac{\Gamma \vdash p \qquad \Gamma \vdash q}{\Gamma \vdash p \land q}
\ (\land\text{-intro}),
\]
and the correspondence is exact: \texttt{constructor} applies the
rule, and the two \texttt{exact} calls discharge the rule's two
premises — the same two goals the \texttt{·} bullets separate above.
\end{experiment}

\section{Elimination Rules: \texttt{apply}, \texttt{exact}, \texttt{cases}}

\texttt{exact} closes a goal outright by supplying a term whose type
matches it precisely — the tactic-mode spelling of simply writing
the term, as every proof in Chapter~\ref{ch:propositions} did.
\texttt{apply} is more interesting: given a hypothesis
$h : p \to q$ and a goal $q$, \texttt{apply h} replaces the goal $q$
with the new goal $p$, since producing a proof of $p$ is now
sufficient (compose it with $h$) to produce the original proof of
$q$. This is $\to$-elimination — modus ponens — read backwards, from
conclusion to premise, which is exactly how a tactic proof is built:
goals shrink toward hypotheses rather than hypotheses growing toward
a conclusion. \texttt{cases} discharges a goal \emph{about} a
hypothesis built from a multi-constructor type by splitting into one
new goal per constructor, each carrying the data that constructor
supplies — for $h : p \lor q$ this yields two goals, one assuming
$p$ and one assuming $q$ ($\lor$-elimination); for $h : \exists x, P\,x$
it yields one goal with a fresh witness variable and a proof that
the witness satisfies $P$ ($\exists$-elimination, taken up properly
once quantifiers themselves are, in Chapter~\ref{ch:quantifiers}).

\begin{experiment}{Modus ponens as \texttt{apply}}
\begin{lean}
theorem chain (p q r : Prop)
    (hpq : p → q) (hqr : q → r) (hp : p) : r := by
  apply hqr
  apply hpq
  exact hp
\end{lean}
Read the goal state as it shrinks: the initial goal is $r$;
\texttt{apply hqr} uses $\Gamma \vdash q \to r$ to reduce it to $q$;
\texttt{apply hpq} uses $\Gamma \vdash p \to q$ to reduce that to
$p$; \texttt{exact hp} closes it. Each \texttt{apply} is one instance
of
\[
\dfrac{\Gamma \vdash p \to q \qquad \Gamma \vdash p}{\Gamma \vdash q}
\ (\to\text{-elim})
\]
read from the conclusion side.
\end{experiment}

\begin{experiment}{Disjunction elimination via \texttt{cases}}
\begin{lean}
theorem or_idem (p : Prop) (h : p ∨ p) : p := by
  cases h with
  | inl hp => exact hp
  | inr hp => exact hp
\end{lean}
The rule being enacted is
\[
\dfrac{\Gamma \vdash p \lor q \qquad \Gamma, p \vdash r
  \qquad \Gamma, q \vdash r}{\Gamma \vdash r}
\ (\lor\text{-elim}),
\]
here with $q$ and $r$ both instantiated to $p$ itself: whichever
disjunct \texttt{cases} produces, the same proof \texttt{hp} closes
the goal, which is exactly why $p \lor p$ collapses to $p$.
\end{experiment}

\section{Exercises}

\begin{exercise}
Prove $(p \land q) \to (q \land p)$ using \texttt{constructor} and
\texttt{exact}.
\end{exercise}

\begin{exercise}
Prove $(p \lor q) \to (q \lor p)$ using \texttt{cases}.
\end{exercise}
